\section{Method}\label{sec:method}

\textbf{Pipeline.} Figure~\ref{fig:arch} shows the work flow. We collected 375 loose ideas and tagged each one with a physics model and a goal period. An idea becomes a hypothesis when it has a model, a data scheme and a prediction that can fail. Each of the 132 hypotheses has a card. The card states the question, the model, the observables, the null and the predictions with kill rules. The card gives a date for each prediction, before any statistic on real data. Claude agents planned and ran the work, one or two hypotheses per agent, under review by the first author.

\textbf{Fields and couplings.} A field is an external drive on all agents. A coupling is an interaction between agents. Four impostors make a field look like a coupling. The scheduler starts and stops agents together. External drives (goals, operator messages, human messages) push all agents at once. Shared model priors make agents of one family write alike. A common input makes agents respond at the same time. Every claim of coupling must remove all four. The strongest test is a partition contrast. A message that the reader read must act more than a message that the reader could not yet see at the same lag.

\textbf{Two layers.} Each hypothesis runs in two layers. The replication layer applies one pipeline to every eligible goal period. The native layer tests the period or natural experiment that fits the hypothesis best. Before real data, each estimator runs on synthetic swarms. These swarms use the real call schedule, room history and roster, with the effect planted or absent. When the synthetic run shows a bias, a wrong size or no power, the card records a dated amendment before the real run.

\textbf{Reserved data.} Before exploration, we reserved 16 goal periods and four natural-experiment windows. No result in this paper uses them. Each hypothesis has a confirmation script, frozen by a hash. The script reads the reserved data only after sign-off by the first author. Three early hypotheses ran on the reserved data (H02, H04, H05). All other results are exploratory.

\textbf{Scoring.} Claude judges every hypothesis (Fig.~\ref{fig:scoring}). First, Claude states the claim that stands in one scoped sentence. Then two Claude judges score the claim. Judge 1 is the coordinator, who saw the whole project. Judge 2 is blind and sees only the card and the evidence. Credence (faithfulness) starts at 0.4, 0.3 or 0.2, by the time when the prediction was written. Each passed check multiplies the odds of the claim. The checks are: the strongest null beaten, replication across periods, survival of data fixes, synthetic recovery, a natural experiment, a rival model beaten, and agreement with known structure. Weak inputs divide the odds. The credence cannot exceed 0.95. Estimated usefulness is the credence times a value if true from 0 to 5. The value depends on the decision that the claim changes. The two scores are averaged in log-odds. A third Claude adjudication settles differences larger than 0.25 in credence, under written rules. A mechanism depth and a fragile flag complete the score.

\textbf{The table.} Each hypothesis names a primary physics model. It can also fit secondary, rival and null models. Figure~\ref{fig:tablezoom} and Appendix~\ref{app:matrix} show the table. Of the 132 claims, the median credence is 0.64, 84 claims reach 0.6, and 4 claims reach mechanism depth M2 (intervention, rival and synthetic recovery).
